% Generated draft: requires mathematical and source review.
\documentclass[11pt]{article}
\usepackage[margin=25mm]{geometry}
\usepackage{fontspec}
\setmainfont{Latin Modern Roman}
\usepackage{amsmath,amssymb,mathtools,unicode-math}
\setmathfont{Latin Modern Math}
\usepackage{enumitem,needspace,xurl,hyperref}
\setlength{\parindent}{0pt}
\setlength{\parskip}{5pt}
\setlength{\emergencystretch}{3em}
\newcommand{\N}{\mathbb{N}}
\newcommand{\Z}{\mathbb{Z}}
\newcommand{\Q}{\mathbb{Q}}
\newcommand{\R}{\mathbb{R}}
\newcommand{\C}{\mathbb{C}}
\newcommand{\F}{\mathbb{F}}
\newcommand{\A}{\mathbb{A}}
\newcommand{\PP}{\mathbb{P}}
\newcommand{\E}{\mathbb{E}}
\newcommand{\eps}{\varepsilon}
\newcommand{\dd}{\,\mathrm{d}}
\begin{document}
\section*{UnsolvedMath ID: 30005485}
\textbf{Source problem code:} catalogue entry
\subsection{Definitions and mathematical statement}
Use the standard convention that a graphon is a symmetric Lebesgue-measurable function \(W\colon[0,1]^2\to[0,1]\). For each integer \(n\geq 2\), construct \(G(n,W)\) on vertex set \(\{1,\ldots,n\}\) by sampling independent uniform random variables \(X_1,\ldots,X_n\) on \([0,1]\), and then, conditional on these variables, including each unordered pair \(\{i,j\}\), \(i<j\), independently with probability \(W(X_i,X_j)\). There are no loops. For a finite graph \(H\), its chromatic number \(\chi(H)\) is the least number of colors in a vertex coloring in which adjacent vertices receive different colors. Expectation below is over both the sampled variables and the conditional edge choices, and \(\log\) denotes the natural logarithm.

The supplied question asks whether, for every fixed graphon \(W\), the sequence
\[
 a_n(W)=\frac{\log n}{n}\,\mathbb{E}\bigl[\chi(G(n,W))\bigr]
\]
converges as \(n\to\infty\).

\textbf{Unresolved qualification.} The supplied record does not specify whether convergence must have a finite limit or may include convergence to \(+\infty\). The two possible targets are
\[
 \forall W\;\exists L_W\in[0,\infty)\quad a_n(W)\longrightarrow L_W
\]
and
\[
 \forall W\;\exists L_W\in[0,\infty]\quad a_n(W)\longrightarrow L_W,
\]
where the second display permits an infinite limit. These targets are distinct. No restriction such as an upper bound strictly below \(1\) on \(W\) is supplied. Accordingly, this draft preserves the convergence question while leaving its intended limit convention and any omitted graphon hypotheses for source review. The graphon and sampling definitions above explicate the standard meaning of the supplied notation; the record itself does not spell them out.
\subsection{Short English statement}
For each fixed graphon \(W\), form its random graph on \(n\) sampled vertices. Does its expected chromatic number, multiplied by \((\log n)/n\), converge as \(n\) grows? The supplied record leaves unclear whether an infinite limit is allowed or further restrictions on \(W\) are intended.
\subsection{Sources}
\begin{itemize}
\item Supplied local UnsolvedMath registry record, permanent ID \(30005485\), title \emph{Convergence of Expected Graphon Chromatic Numbers}, exact source problem code \texttt{catalogue entry}, in the collection \emph{Oberwolfach Reports Open Problems}. The record gives proposed year 2023 and names no proposer for the graphon question.
\item Supplied report citation: \emph{Random Graphs: Combinatorics, Complex Networks and Disordered Systems} (2023), Oberwolfach Reports, pp.\ 883--914, DOI \texttt{10.4171/owr/2023/16}. \url{https://doi.org/10.4171/owr/2023/16}.
\item Supplied literature reference: \emph{Graphon chromatic polynomials}. \url{https://arxiv.org/abs/2109.07773}.
\item Supplied literature reference: \emph{OWR 2023 report}. \url{https://ems.press/content/serial-article-files/47010?nt=1}.
\item Literature assessment dated exactly \texttt{2026-08-22}: the supplied curation labels the problem \texttt{partially\_solved} and states that convergence is proved for several graphon classes and general bounds are known, while the convergence question remains unresolved in full generality. This is supplied dated curation, not an independently verified literature assessment; its status and cited claims require independent verification.
\item The supplied statement status is \texttt{corrected\_verified}; its verification field reports comparison with the official Oberwolfach report. This draft records that supplied verification claim without claiming to have checked the report. The raw extraction also contains a separate hypergraph-matching passage attributed to Matthew Kwan, referring to a conjecture with Asaf Ferber; that attribution is not supplied for the graphon question.
\end{itemize}
\end{document}
